\bookchapter{Functions and graphs}\label{ch:functions}

\section*{Explanation}
A \emph{function} is a rule that assigns to every allowed input exactly one output. We write $y=f(x)$ and read it as ``$y$ equals $f$ of $x$''. The set of allowed inputs is the \emph{domain}, and the set of outputs that actually occur is the \emph{range}.

\begin{center}
\begin{tikzpicture}[>=Stealth, box/.style={draw,thick,rounded corners,minimum height=1.3cm,minimum width=3.4cm,fill=orange!15}]
  \node (in) at (0,0) {input $x=3$};
  \node[box] (f) at (4.8,0) {$f(x)=2x+1$};
  \node (out) at (9.6,0) {output $f(3)=7$};
  \draw[->,thick] (in) -- (f);
  \draw[->,thick] (f) -- (out);
\end{tikzpicture}

\captionof{figure}{A function as a machine: every input produces exactly one output.}\label{fig:b1_1}
\end{center}

The \emph{graph} of $f$ is the set of all points $(x,f(x))$ in the coordinate plane. We can draw it by computing a table of values, plotting the points and joining them with a smooth curve. For $f(x)=x^2-1$:
\begin{center}
\renewcommand{\arraystretch}{1.3}
\begin{tabular}{c|ccccc}
$x$ & $-2$ & $-1$ & $0$ & $1$ & $2$\\\hline
$f(x)$ & $3$ & $0$ & $-1$ & $0$ & $3$
\end{tabular}
\end{center}

\begin{center}
\begin{tikzpicture}
\begin{axis}[gr, xmin=-3,xmax=3,ymin=-2,ymax=5, xtick={-3,...,3}, ytick={-2,...,5}]
  \addplot[blue,domain=-2.4:2.4]{x^2-1};
  \addplot[only marks,mark=*,red,mark size=2.5pt] coordinates {(-2,3)(-1,0)(0,-1)(1,0)(2,3)};
\end{axis}
\end{tikzpicture}

\captionof{figure}{The graph of $f(x)=x^2-1$. The red points come from the table.}\label{fig:b1_2}
\end{center}

\textbf{Vertical line test.} A curve is the graph of a function exactly when every vertical line crosses it at most once. A circle is therefore \emph{not} the graph of a function.

\begin{example}
For $f(x)=x^2-1$ we get $f(3)=8$ and $f(-2)=3$. The domain is all real numbers, the range is $[-1,\infty)$ because the smallest value is $f(0)=-1$.
\end{example}

\begin{example}
Domains: $\sqrt{x}$ needs $x\ge0$, and $\dfrac1x$ needs $x\neq0$.
\end{example}

\section*{Linear functions}
A \emph{linear function} has the form $f(x)=mx+b$. Its graph is a straight line. The number $m$ is the \emph{slope} (how much $y$ rises when $x$ grows by $1$), and $b$ is the \emph{$y$-intercept} (the value at $x=0$). From two points on the line we get
\[
m=\frac{y_2-y_1}{x_2-x_1}=\frac{\text{rise}}{\text{run}}.
\]

\begin{center}
\begin{tikzpicture}
\begin{axis}[gr, width=8.5cm,height=7cm, xmin=-3,xmax=4,ymin=-4,ymax=7, xtick={-3,...,4}, ytick={-4,...,7}, legend pos=south east]
  \addplot[blue,domain=-3:4]{2*x+1};
  \addlegendentry{$y=2x+1$}
  \addplot[red,domain=-3:4]{3-x};
  \addlegendentry{$y=-x+3$}
  \addplot[green!50!black,domain=-3:4]{0.5*x-2};
  \addlegendentry{$y=\tfrac12x-2$}
  \draw[blue!70!black,dashed,very thick] (axis cs:0,1) -- (axis cs:1,1) node[midway,below,font=\scriptsize]{run $=1$} -- (axis cs:1,3) node[midway,right,font=\scriptsize]{rise $=2$};
\end{axis}
\end{tikzpicture}

\captionof{figure}{Three lines. A larger slope means a steeper line, a negative slope means the line falls. The dashed triangle shows rise and run for the blue line.}\label{fig:b1_3}
\end{center}

\begin{example}
The line through $(1,3)$ and $(3,7)$ has slope $m=\frac{7-3}{3-1}=2$. Using $y=mx+b$ at the point $(1,3)$ gives $3=2\cdot1+b$, so $b=1$ and the line is $y=2x+1$.
\end{example}

\subsection*{How to draw a line by hand}
You do not need a table of many values to draw a line. A line is fixed by \emph{one point and the slope}, and two points are enough to place a ruler. There are two good ways to start.

\begin{keyidea}[Drawing a line in four steps]
\begin{enumerate}[nosep,label=\textbf{\arabic*.}]
  \item \textbf{Find a starting point.} Take the $y$-intercept $(0,b)$, or the $x$-intercept: set $y=0$ and solve $mx+b=0$, which gives $x=-\frac{b}{m}$.
  \item \textbf{Write the slope as a fraction} $m=\frac{\text{rise}}{\text{run}}$. For example $2=\frac{2}{1}$, and $-\frac34=\frac{-3}{4}$.
  \item \textbf{Walk from the starting point:} go \emph{run} steps to the right, then \emph{rise} steps up. If the rise is negative, go down instead.
  \item \textbf{Mark the new point and draw the line} through both points with a ruler. Extend it past the points in both directions.
\end{enumerate}
\end{keyidea}

\textbf{Way 1: start at the $x$-intercept.} The $x$-intercept is the point where the line crosses the $x$-axis, i.e.\ where $y=0$.

\begin{example}
Draw $y=2x-4$. Set $y=0$: $0=2x-4$, so $x=2$. The line crosses the $x$-axis at $(2,0)$. The slope is $2=\frac21$: from $(2,0)$ go $1$ to the right and $2$ up to reach $(3,2)$. Draw the line through $(2,0)$ and $(3,2)$ (left panel of the figure below). It also passes through $(0,-4)$, the $y$-intercept.
\end{example}

\textbf{Way 2: start at the $y$-intercept.} This is often even quicker, because $b$ can be read directly from the equation.

\begin{example}
Draw $y=-\frac34x+3$. The $y$-intercept is $(0,3)$. The slope is $-\frac34=\frac{-3}{4}$: from $(0,3)$ go $4$ to the right and $3$ down to reach $(4,0)$. Draw the line through $(0,3)$ and $(4,0)$ (right panel of the figure).
\end{example}

\begin{center}
\begin{tikzpicture}
\begin{axis}[gr, width=6.6cm,height=6.2cm, anchor=south west, at={(0,0)}, xmin=-1,xmax=5,ymin=-5,ymax=5, xtick={-1,...,5}, ytick={-5,...,5}, title={\small $y=2x-4$}]
  \addplot[blue,domain=-1:4.5]{2*x-4};
  \addplot[only marks,mark=*,mark size=2pt,red] coordinates {(2,0) (3,2)};
  \addplot[only marks,mark=*,mark size=2pt,black] coordinates {(0,-4)};
  \draw[red!70!black,dashed,very thick,-{Stealth}] (axis cs:2,0) -- (axis cs:3,0) node[midway,below,yshift=-9pt,font=\scriptsize]{run 1};
  \draw[red!70!black,dashed,very thick,-{Stealth}] (axis cs:3,0) -- (axis cs:3,2) node[midway,right,font=\scriptsize]{rise 2};
  \node[font=\scriptsize,anchor=west] at (axis cs:0.15,-4) {$y$-intercept};
\end{axis}
\begin{axis}[gr, width=6.6cm,height=6.2cm, anchor=south west, at={(7.6cm,0)}, xmin=-1,xmax=6,ymin=-2,ymax=5, xtick={-1,...,6}, ytick={-2,...,5}, title={\small $y=-\tfrac34x+3$}]
  \addplot[blue,domain=-1:5.5]{-0.75*x+3};
  \addplot[only marks,mark=*,mark size=2pt,red] coordinates {(0,3) (4,0)};
  \draw[red!70!black,dashed,very thick,-{Stealth}] (axis cs:0,3) -- (axis cs:4,3) node[midway,above,font=\scriptsize]{run 4};
  \draw[red!70!black,dashed,very thick,-{Stealth}] (axis cs:4,3) -- (axis cs:4,0) node[midway,right,font=\scriptsize]{3 down};
\end{axis}
\end{tikzpicture}

\small\textbf{Drawing lines from the intercepts.} Left: start at the $x$-intercept $(2,0)$, go $1$ right and $2$ up. Right: start at the $y$-intercept $(0,3)$, go $4$ right and $3$ down.
\end{center}

\textbf{Check your line.} Pick any $x$ that you did not use, for example $x=1$ in the first example: $y=2\cdot1-4=-2$. The point $(1,-2)$ must lie on your drawn line. If it does not, one of the steps went wrong.

\textbf{Typical mistakes.}
\begin{itemize}[nosep]
  \item Mixing up rise and run: the \emph{run} (left-right) always comes first, then the \emph{rise} (up-down).
  \item Forgetting the sign: a negative slope means the line goes \emph{down} when you move to the right.
  \item Stopping at the marked points: the line goes on in both directions, so extend it.
\end{itemize}

\section*{Quadratic functions}
A \emph{quadratic function} $f(x)=ax^2+bx+c$ (with $a\neq0$) has a parabola as its graph. If $a>0$ it opens upward, if $a<0$ downward. The turning point is the \emph{vertex}, located at $x=-\dfrac{b}{2a}$. The points where the graph meets the $x$-axis are the \emph{roots}, the solutions of $f(x)=0$.

\begin{center}
\begin{tikzpicture}
\begin{axis}[gr, width=9cm,height=7cm, xmin=-3.5,xmax=5.5,ymin=-4,ymax=9, xtick={-3,...,5}, ytick={-4,-2,...,8}, legend pos=north east]
  \addplot[blue,domain=-3.5:5.5]{x^2};
  \addlegendentry{$y=x^2$}
  \addplot[red,domain=-3.5:5.5]{4-x^2};
  \addlegendentry{$y=-x^2+4$}
  \addplot[green!50!black,domain=-3.5:5.5]{(x-2)^2-1};
  \addlegendentry{$y=(x-2)^2-1$}
  \addplot[only marks,mark=*,blue] coordinates {(0,0)};
  \addplot[only marks,mark=*,red] coordinates {(0,4)};
  \addplot[only marks,mark=*,green!50!black] coordinates {(2,-1)};
\end{axis}
\end{tikzpicture}

\captionof{figure}{Three parabolas with their vertices marked. The term $(x-2)$ shifts the graph $2$ to the right, and the $-1$ shifts it $1$ down.}\label{fig:b1_4}
\end{center}

\begin{example}
$f(x)=x^2-4x+3$. The vertex is at $x=\frac{4}{2}=2$ with $f(2)=4-8+3=-1$. Since $x^2-4x+3=(x-1)(x-3)$ the roots are $x=1$ and $x=3$.
\end{example}

\section*{Exponential and logarithmic functions}
An \emph{exponential function} $f(x)=a^x$ with $a>0$ grows very fast when $a>1$ and decays when $0<a<1$. It always passes through $(0,1)$ and its values are always positive. A very important base is Euler's number $e\approx2.71828$, and $e^x$ is also written $\exp(x)$.

The \emph{logarithm} undoes the exponential: $\log_a(y)=x$ means $a^x=y$. The logarithm with base $e$ is the \emph{natural logarithm} $\ln$. It is defined only for $x>0$ and passes through $(1,0)$. The graph of $\ln x$ is the graph of $e^x$ mirrored at the line $y=x$.

\begin{center}
\begin{minipage}{0.49\textwidth}
\centering
\begin{tikzpicture}
\begin{axis}[gr, width=7cm,height=6cm, xmin=-3,xmax=3,ymin=-1,ymax=8, xtick={-3,...,3}, ytick={-1,...,8}, legend pos=north west]
  \addplot[blue,domain=-3:3]{2^x};
  \addlegendentry{$2^x$}
  \addplot[red,domain=-3:2.1]{exp(x)};
  \addlegendentry{$e^x$}
  \addplot[green!50!black,domain=-3:3]{0.5^x};
  \addlegendentry{$(\tfrac12)^x$}
  \addplot[only marks,mark=*,black] coordinates {(0,1)};
\end{axis}
\end{tikzpicture}
\end{minipage}\hfill
\begin{minipage}{0.49\textwidth}
\centering
\begin{tikzpicture}
\begin{axis}[gr, width=7cm,height=7cm, axis equal image, xmin=-2,xmax=5,ymin=-2,ymax=5, xtick={-2,...,5}, ytick={-2,...,5}, legend pos=south east]
  \addplot[red,domain=-2:1.6]{exp(x)};
  \addlegendentry{$e^x$}
  \addplot[blue,domain=0.14:5]{ln(x)};
  \addlegendentry{$\ln x$}
  \addplot[black,dashed,domain=-2:5,thin]{x};
  \addlegendentry{$y=x$}
\end{axis}
\end{tikzpicture}
\end{minipage}

\captionof{figure}{Left: exponential functions all pass through $(0,1)$. Right: $e^x$ and $\ln x$ are mirror images of each other at the dashed line $y=x$.}\label{fig:b1_5}
\end{center}

\section*{Composition of functions}
We can feed the output of one function into another. The \emph{composition} is $(f\circ g)(x)=f(g(x))$: first $g$, then $f$.

\begin{example}
With $f(x)=2x+1$ and $g(x)=x^2$: $f(g(3))=f(9)=19$, but $g(f(3))=g(7)=49$. The order matters.
\end{example}

\begin{aibox}
A single neuron computes $a(x)=\sigma(wx+b)$. This is a linear function $wx+b$ followed by an activation function $\sigma$, i.e.\ a composition. A neural network is a long composition of such functions, so everything in a network is built from the ideas of this chapter.
\end{aibox}

\section*{Practice problems}
\begin{enumerate}[label=\textbf{\thechapter.\arabic*}]
  \item For $f(x)=3x-2$ compute $f(0)$, $f(2)$ and $f(-1)$.
  \item For $g(x)=x^2-4x+1$ compute $g(3)$ and $g(-1)$.
  \item Give the slope and the $y$-intercept of the line $2y-6x=8$.
  \item Find the linear function through the points $(1,3)$ and $(3,11)$.
  \item Describe the graph of $y=(x-3)^2+2$: in which direction does it open, where is the vertex, and why does it have no root?
  \item Evaluate without a calculator: $2^3$, $2^{-2}$, $\log_2 8$, $\ln(e^5)$, $\log_{10}0.01$.
  \item Let $f(x)=2x+1$ and $g(x)=x^2$. Compute $f(g(2))$ and $g(f(2))$.
  \item Name two points on the graph of $y=2^x$. Use them to give two points on the graph of $y=\log_2 x$.
  \item $\star$ For $f(x)=x^2-6x+5$ find the vertex, the roots and the $y$-intercept, and sketch the graph.
\end{enumerate}

% =====================================================
